%%%%%%%%%%%%%%%%%%%%%%%%%%%%%%%%%%%%%%%%%%%%%%%%%%%%%
%			HLAVIČKA								%
%%%%%%%%%%%%%%%%%%%%%%%%%%%%%%%%%%%%%%%%%%%%%%%%%%%%%
\documentclass[openany,12pt]{memoir}

\usepackage[utf8]{inputenc} 
\usepackage[czech]{babel}
\usepackage[T1]{fontenc}
\usepackage[top=1.5cm, bottom=2cm, left=2cm, right=2cm]{geometry}  % --> NASTAVENÍ OKRAJŮ
\usepackage{fancyhdr}
\usepackage{graphicx}
%\usepackage{lmodern}
\usepackage{tgschola}  % --- Silnější font
\usepackage{xwatermark}
\usepackage{xcolor}
\usepackage{changepage}
\usepackage{newunicodechar}



%%%%%% Package na zpěvník
\usepackage[full]{leadsheets}%http://mirrors.nic.cz/tex-archive/macros/latex/contrib/leadsheets/leadsheets_en.pdf   --> dokumentace	
\definesongtitletemplate{empty}{} 
\setchords{
format = \bfseries,   %tučné akordy
minor = {mi},% 
input-notation = {german},%
output-notation = {german}%
}
\definesongtitletemplate{empty}{} 

\newlength{\drop}
\newwatermark[allpages,color=red!50,angle=0,scale=2, xpos=0,ypos=0]{\includegraphics[width=5cm]{obr/pozadi2.jpg}} %--> dvojka na pozadí


%%%% Vlastní příkazy
\newcounter{Slokočet}   %Automatické číslování slok
\newcommand{\mezera}{\vspace*{0.5cm}}   %Horizontální odsazení slok
\newcommand{\stred}{5.2cm}   %%% Na zarovnání slok doprostřed, pozn. automatičtější zarovnávání na střed nejde
\newcommand{\refren}{\mezera \noindent \textbf{R:} } %refrén
\newcommand{\sloka}{\mezera \noindent \addtocounter{Slokočet}{1} \arabic{Slokočet}. } 	%sloka, která se automaticky čísluje
\newcommand{\ssloka}{\mezera \noindent}  % vlastní číslo sloky
\newcommand{\carka}{,\:}
\newcommand{\predehra}{\mezera \hspace*{-1.85cm} \textbf{Předehra:} }
\newcommand{\m}[1]{\color{white}{#1}}  %Pro akordy
\newcommand{\ap}{'} %Pro apostrof
\newcommand{\elipsa}{\kern\fontdimen3\font} %Příkaz pro lepší zacházení s výpustkami (=...); je to vpodstatě jen mezera mezi tečkama výpustky

%%%%%%%%%%%%%%%%%%%%%%%%%%%%%%%%%%%%%%%%%%%%%%%%%%%%%%
%			NÁVOD									 %
%%%%%%%%%%%%%%%%%%%%%%%%%%%%%%%%%%%%%%%%%%%%%%%%%%%%%%
%1. Věci v hlavičce IGNOROVAT
%2. Píseň psát do prostoru mezi \begin{song} a \end{song}
%3. další řádek se značí dvěma odsazeníma (= dvakrát stisknout enter)
%4. \refren vždy na začátku refrenu a \sloka na začátek sloky (automaticky se čísluje)
%  \ = alt gr + q ; [/] = alt gr f/g ; {/} = alt gr + b/n; ^ = alt gr + 3 + mezera
%Cokoliv napíšete do ^{  } se bude brát jako akord
%když se toto bude dotýkat nějakého slova (nebude mezi tím a slovem mezera)
%tak se akordy zjeví nad slovem, ale pište to před slovo
%Když se to nedotýká slova, tak akord lítá ve vzduchu a vytiskne se větší mezera
%První možnost je asi preferovanější
%5. Akordy stačí psát jen do první sloky, když se nezmění -- kytaristi to zvládnou



%\usepackage{subfiles}
%</++>

\begin{document}

\pagestyle{empty}
%\addtocontents{toc}{\protect\thispagestyle{empty}}
%\tableofcontents \thispagestyle{empty}\newpage
\newgeometry{top=1.5cm, bottom = 0cm, left = 2cm, right = 2cm}

\newpage

%Pro kompilaci je důležitý, 	aby ten input a \end{document} zůstali na posledních
%dvou řádcích!!!!
\begin{song}{title=\predtitle\centering Kaliforňané \\\large Greenhorns  \vspace*{-0.3cm}}  %% sem se napíše jméno songu a autor
\begin{centerjustified}
\nejnejvetsi

\sloka
	Ať ^{C\z}všichni ^{Ami}Kaliforňané se ^{G}za nás ^{\z C}pomodlí,

	kdo táhneš s mulou ^{F\z}přes pláně, ^{\z C}rozluč se s ^{\z F}pohodlím,

   	^{C\z }musíš si taky ^{F\z}nakoupit ^{C}solený ^{\z F}vepřový,

   	a ^{C}kdo má štěstí, ^{Ami\z}po~cestě i ^{G\z}něco ^{\z C}uloví.

\sloka
	Po boku měj dva kolťáky a vostrej Bowie nůž.

   	Jestli tu cestu překonáš, potom seš celej muž.

   	Zelenáči, až uvidíš kol sebe poušť a prach

   	a nikde živý stvoření, dostaneš asi strach.

\sloka
	Dva měsíce ti potrvá ta cesta přes pláně,

   	koukej, ať vždycky z čeho máš vařit svý snídaně,

   	protože k čertu dojedeš při bídným vaření,

   	i tvoje mula musí žrát, jináč ti zcepení.

 \sloka
	A taky proti rudochům drž stráže za noci,

   	po dlouhejch denních pochodech propuknou nemoci,

   	horečka, to je špatná věc, ta lidi poráží,

   	a kdo je marod, nemůže stát v noci na stráži.

 \sloka
	Břitvu nech doma, do tváře se radši nedívej,

   	až přijdeš, budeš zarostlej a taky šedivej,

   	až tě tvá milá uvidí, uslyšíš asi pláč,

   	svý mámě řekneš: ,,Podívej, to je tvůj zelenáč!{}``

\sloka
	U Plat je nouze vo dřevo, když tě to nezdolá,

   	omáčku si tam zamícháš paznehtem z buvola,

   	u Plat je nouze vo klacek, zem je tam jalová,

   	i já bych teď už radši byl u Leopoldova. %mysu Hornova 

\end{centerjustified}
\setcounter{Slokočet}{0}
\end{song}

\end{document}
